\section{Identification Strategy}
\label{sec:identification}

We identify the causal effect of a credit-risk disclosure on depositor
behavior.  Cross-sectional comparisons conflate information responses
with pre-existing sorting \citep{covitz2004market, chen2022bank}: riskier
banks carry fewer uninsured deposits in equilibrium.  We exploit the
CECL Day-One adjustment as a natural experiment---timing set by
regulation, and magnitude reflecting the bank's assessment of lifetime
expected losses on its existing portfolio.  The cross-sectional intensity
of the shock is therefore largely fixed by past underwriting rather than
by contemporaneous deposit dynamics, a claim we treat as testable and
probe with extended pre-trends, residualized treatments, and the
within-bank composition design below.

\subsection{Empirical Design}
\label{sec:identification:did}

Our baseline is a difference-in-differences estimator:
\begin{equation}
  y_{it}
  = \alpha_i + \gamma_t
  + \beta \cdot \mathrm{Post}_{it} \times \mathrm{CECL}_i
  + \mathbf{X}_{it-1}'\boldsymbol{\delta}
  + \varepsilon_{it},
  \label{eq:did}
\end{equation}
where $y_{it}$ is the outcome for bank $i$ in quarter $t$; $\alpha_i$
and $\gamma_t$ are bank and calendar-quarter fixed effects;
$\mathrm{Post}_{it}$ equals one at or after bank $i$'s adoption date;
$\mathrm{CECL}_i$ is either continuous \texttt{CECL Adj} or binary
\texttt{High CECL}; and $\mathbf{X}_{it-1}$ collects lagged controls
(log assets, equity-to-assets, interest expense-to-assets).  Bank fixed
effects absorb all time-invariant differences; quarter fixed effects
absorb aggregate macroeconomic forces common to all banks in a quarter,
including the tightening cycle and any nationwide reaction to the 2023
banking stress.  We also estimate an event-study specification with
event-time dummies $k\in[-10,10]$, $k=-1$ omitted; the baseline window
excludes 2020Q1--Q2 but we re-estimate the full 2018Q1-onward window in
Section~\ref{sec:robustness:pretrends}.  Standard errors are clustered at
the bank level throughout.\footnote{Because 92 percent of the sample
adopts CECL in 2023Q1, the design has a strong common-shock component,
and inference clustered only at the bank level could understate
uncertainty if cross-sectional residuals share the post-shock
disturbance.  Two-way clustering by bank and calendar quarter typically
inflates standard errors modestly under this design, but the sign and
rank ordering of the reported coefficients are unchanged under coarser
inference; we flag two-way and design-based inference for staggered
adoption as areas for further work, and view the qualitative
conclusions of the paper as robust to the specific clustering choice.}

\subsection{Treatment Exogeneity}
\label{sec:identification:predetermination}

The 2023Q1 adoption date was set by regulation and cannot have been
chosen to coincide with favorable deposit conditions.  The size of the
Day-One adjustment reflects the cumulative composition of the loan
portfolio---long maturities, weak borrower credit, thin collateral---%
fixed by prior underwriting.
Table~\ref{tab:cecl_predictors_cross_section} confirms the adjustment
loads on NPL ratios, size, and capital in the expected directions.
Adjusted $R^2$ maxes at roughly 1.4 percent in an expanded specification
adding deposit composition and pre-adoption performance trends.  Low
explanatory power does not prove exogeneity, but it establishes that
high-CECL banks are not simply the observably weak, deposit-fragile, or
stress-exposed banks relabeled.  We lean on three complementary pieces
of evidence: flat extended pre-trends, robustness to a treatment
residualized on the full expanded predictor set, and the within-bank
composition design of the triple difference.  The standard was public
from June 2016, but the \emph{magnitude} of any given bank's Day-One
charge was not knowable in advance; anticipation of the event date
without knowledge of intensity does not bias the DiD, and anticipation
of intensity would appear as pre-adoption divergence.

Two components of the Day-One magnitude deserve separate treatment.
The first is legacy loan composition---maturity structure, borrower
credit quality, collateral coverage---which was fixed by prior
underwriting and is uncorrelated with 2023 deposit dynamics under any
reasonable assumption.  The second is the forecast and modeling
component required by ASC~326: banks combine historical loss experience
with reasonable-and-supportable forward-looking forecasts of
macroeconomic conditions and choose a lifetime-loss estimation approach.
This component is a contemporaneous choice at adoption and is in
principle correlated with 2022--2023 macro conditions.  Our residualized
treatment exercise
(Section~\ref{sec:robustness:resid}) is designed for exactly this
concern: we strip the Day-One adjustment of everything predictable from
the full 2022Q4 observable set, including deposit composition and
eight-quarter performance trends, and re-estimate with the residual.
The estimate attenuates but retains its sign at roughly two-thirds the
baseline magnitude, consistent with a bound on the modeling-choice
channel rather than proof of exogeneity.  We are careful throughout not
to claim more than the design delivers: the Day-One adjustment is
predetermined with respect to legacy loan composition, and any
remaining forward-looking modeling choice enters as a noise-plus-signal
component that the residualization partially isolates.

\subsection{Information Content}
\label{sec:identification:mechanism}

Prior depositor-discipline tests typically exploit realized distress in
which new information arrives simultaneously with deterioration in cash
flows or liquidity \citep{flannery1996evidence, martinez2001depositors,
iyer2016tale}.  The CECL Day-One adjustment separates the two: it
records a pure accounting charge that does not trigger a cash transfer,
alter borrower repayment terms, reduce bank liquidity, or change the
risk of the underlying loans.  A depositor who withdraws or demands a
higher rate responds to newly disclosed credit-risk information, not to
concurrent worsening of the bank's actual condition.

\subsection{The 2023 Banking Stress}
\label{sec:identification:stress}

The principal threat is the coincidence of the primary adoption quarter
with the March 2023 bank failures.  Calendar-quarter fixed effects
absorb only the \emph{common} component of that shock; the sharper
concern is that exposure to the 2023 stress varied across banks with
characteristics that correlate with the CECL adjustment
\citep{caglio2024depositors, jiang2024monetary}.
Section~\ref{sec:robustness:stress} confronts this with
characteristics-by-quarter controls and staggered-cohort event
studies.\footnote{As a supplementary check, event studies on the large
publicly traded banks that adopted CECL in 2020Q1 show flat market-cap
and CDS-spread pre-trends and sharp adjustments at adoption (results
available on request); the pandemic's simultaneous arrival makes this
evidence suggestive rather than definitive, and our community-bank tests
do not rely on this cohort.}
